\documentclass[a4paper,12pt]{ctexart}
\usepackage{xcolor}
\usepackage{graphicx}
\usepackage{amsmath}
\usepackage[top=1.8cm, bottom=1.8cm, left=1.8cm, right=1.8cm]{geometry}
\usepackage{array}
\usepackage{hyperref}
\hypersetup{colorlinks=true,linkcolor=blue!70!black}
\linespread{1.35}

\definecolor{defblue}{RGB}{0,80,160}
\definecolor{thinkorange}{RGB}{230,120,20}
\definecolor{formulared}{RGB}{180,30,50}
\definecolor{funboxbg}{RGB}{245,248,252}
\definecolor{examplebg}{RGB}{255,250,240}
\definecolor{practicegreen}{RGB}{40,130,70}

\newenvironment{thinkbox}{
  \vspace{0.3cm}
  \begin{center}
  \begin{minipage}{0.88\textwidth}
  \colorbox{funboxbg}{
  \begin{minipage}{0.94\textwidth}
  \vspace{0.15cm}
  \textbf{\textcolor{thinkorange}{【 想一想 】}}
}{
  \vspace{0.15cm}
  \end{minipage}}
  \end{minipage}
  \end{center}
  \vspace{0.3cm}
}

\newenvironment{definition}[1]{
  \vspace{0.2cm}
  \begin{center}
  \begin{minipage}{0.88\textwidth}
  \fbox{
  \begin{minipage}{0.92\textwidth}
  \vspace{0.1cm}
  \textbf{\textcolor{defblue}{【 #1 】}}
}{
  \vspace{0.1cm}
  \end{minipage}}
  \end{minipage}
  \end{center}
  \vspace{0.2cm}
}

\newenvironment{example}{
  \vspace{0.2cm}
  \begin{center}
  \begin{minipage}{0.88\textwidth}
  \colorbox{examplebg}{
  \begin{minipage}{0.94\textwidth}
  \vspace{0.15cm}
  \textbf{\textcolor{orange!80!black}{【 例题 】}}
}{
  \vspace{0.15cm}
  \end{minipage}}
  \end{minipage}
  \end{center}
  \vspace{0.2cm}
}

\newenvironment{practice}{
  \vspace{0.3cm}
  \begin{center}
  \begin{minipage}{0.88\textwidth}
  \colorbox{funboxbg}{
  \begin{minipage}{0.94\textwidth}
  \vspace{0.15cm}
  \textbf{\textcolor{practicegreen}{【 练一练 】}}
}{
  \vspace{0.15cm}
  \end{minipage}}
  \end{minipage}
  \end{center}
  \vspace{0.2cm}
}

\newenvironment{figplaceholder}[1]{
  \vspace{0.2cm}
  \begin{center}
  \fbox{
  \begin{minipage}{0.82\textwidth}
  \centering
  \textcolor{blue!70!black}{\textbf{【插图位置】}} \\[0.2cm]
  #1
  \end{minipage}}
  \end{center}
  \vspace{0.2cm}
}

\newcommand{\formula}[1]{\begin{center}\textcolor{formulared}{\large\boxed{#1}}\end{center}}

\title{\textcolor{blue!70!black}{\Huge\textbf{物理1 · 力学入门}}}
\author{}
\date{}

\begin{document}
\maketitle
\thispagestyle{empty}

\section*{\textcolor{blue!70!black}{致同学}}

欢迎来到物理的世界！如果你已经读过《自然科学小故事》，你一定还记得苹果树下的牛顿、浴缸里的阿基米德、看烧水壶的小瓦特——那些故事让你对"为什么"充满了好奇。

现在，我们要把这些"为什么"整理成一套系统的思维工具。你会学到物理学家是怎么描述运动的、怎么计算力的、怎么用简单的数学把复杂的问题变清晰。在这个过程中，你会遇到\textbf{公式}——但别紧张，它们不是什么可怕的东西。公式其实就是"用数学符号写的一句话"，就像你学会了英语可以说"I am happy"，学会了物理公式就可以说"$v = s/t$"，意思是"速度等于路程除以时间"。

准备好了吗？让我们正式开始。

\newpage

\section{\textcolor{orange!80!black}{第一课\quad 运动的世界——快与慢}}

\subsection*{什么在动？什么没动？}

你坐在教室里，老师在黑板上写字。你觉得——自己没动，老师也没动。但如果有人在操场上透过窗户看进来，他会发现：整个教室都在动——因为地球在自转，你正以每秒约460米的速度（北京纬度）向东飞驰！

所以，说一个物体"在动"还是"没动"，首先要选定一个\textbf{参照物}。参照物是你假定为"不动"的物体，然后看研究对象相对于它是动了还是没动。

\begin{definition}{参照物}
判断一个物体是否运动，需要先选一个假定不动的物体作为\textbf{参照物}。同一个物体，选不同的参照物，得出的运动结论可能不同。
\end{definition}

\textbf{举例：}你坐在地铁上，以车厢为参照物——你没动；以站台上的广告牌为参照物——你在飞快地移动。古代诗句"两岸青山相对出"，诗人以船为参照物，看到的青山就是"移动"的。

\subsection*{速度——运动的"快慢分"}

两个人在操场上跑步，你怎么判断谁跑得快？一样的时间，谁跑得远谁快；一样的距离，谁用时短谁快。物理学家把这两个直觉统一成了一个量——\textbf{速度}。

\begin{definition}{速度}
\textbf{速度}是描述物体运动快慢的物理量，等于路程除以时间。
\end{definition}

\formula{v = \frac{s}{t}}

其中：$v$ 表示速度，$s$ 表示路程，$t$ 表示时间。速度的单位是\textbf{米/秒（m/s）}，生活中更常用的是\textbf{千米/小时（km/h）}。换算关系是：

\formula{1\ \text{m/s} = 3.6\ \text{km/h}}

也就是说，如果一个人每秒跑1米，换算成开车人的说法就是"时速3.6公里"——非常非常慢。

\begin{example}
小明跑了1500米，用了6分钟。他的平均速度是多少？

\textbf{解：} 6分钟 = $6 \times 60 = 360$秒。\\
$v = \dfrac{s}{t} = \dfrac{1500}{360} \approx 4.17\ \text{m/s}$。\\
换算成km/h：$4.17 \times 3.6 \approx 15\ \text{km/h}$。差不多是慢骑自行车的速度。
\end{example}

\subsection*{匀速直线运动}

如果物体沿一条直线运动，并且快慢始终不变——这就是最简单的运动方式：\textbf{匀速直线运动}。

对于匀速直线运动，速度公式可以变形成：

\formula{s = vt \qquad t = \frac{s}{v}}

知道任意两个量，就能算出第三个。物理的魅力就在这里——用简洁的数学找到确定的答案。

在物理学中，我们还常常用\textbf{$s$-$t$图像}来表示运动：横轴是时间 $t$，纵轴是路程 $s$。匀速直线运动的 $s$-$t$ 图像是一条\textbf{倾斜的直线}——倾斜得越厉害，说明速度越大。如果图像是一条水平线，说明路程不随时间变化——物体\textbf{静止}。

\begin{figplaceholder}
此处插入s-t图像示意图：三条不同斜率的直线代表不同速度（见 figure/requirements/physics-mechanics1-ch1.txt 插图1）
\end{figplaceholder}

\subsection*{不是所有运动都是匀速的}

生活里的运动很少是完美的匀速——汽车起步的时候越来越快，刹车的时候越来越慢。这时候我们可以算\textbf{平均速度}：总路程除以总时间。但要注意，平均速度可能和任何一个瞬间的实际速度都不相等——就像你语文数学英语三门课的平均分，可能不等于任何一门的分数。

\begin{thinkbox}
从北京坐高铁到上海，全程约1318公里，大约需要4.5小时。算一下高铁的平均速度是多少km/h？注意这个速度是平均速度——高铁出站和进站的时候实际速度远低于此，而巡航时实际速度可达350 km/h。平均值和瞬时值可以差很多！
\end{thinkbox}

\begin{practice}
1. （判断）小明坐在行驶的公交车上，以车窗为参照物，小明是运动的。（\quad）\\
2. 一辆汽车在平直公路上匀速行驶，10秒内通过了200米，速度是\_\_\_\_\_ m/s，合\_\_\_\_\_ km/h。\\
3. （选择）在匀速直线运动中，下列关于速度公式 $v=s/t$ 的说法正确的是（\quad）\\
\quad A. $v$ 与 $s$ 成正比 \quad B. $v$ 与 $t$ 成反比 \quad C. $v$ 不变时 $s$ 与 $t$ 成正比 \quad D. 以上都不对\\
4. （填空）$s$-$t$ 图像中，一条水平线表示物体处于\_\_\_\_\_状态。\\
5. 从北京到天津约120 km，高铁以300 km/h匀速行驶，需要多少分钟到达？\\
6. （计算）一辆自行车5分钟内通过了1200米，速度是多少m/s？合多少km/h？\\
7. （判断）s-t图像中一条倾斜的直线表示物体在做匀速直线运动。（\quad）\\
8. （选择）下列属于匀速直线运动的是（\quad）\\
\quad A. 绕操场跑步 \quad B. 电梯上升过程 \quad C. 火车在平直轨道上匀速行驶 \quad D. 过山车\\
9. （填空）速度是表示物体\_\_\_\_\_的物理量，公式为\_\_\_\_\_。
\end{practice}

\newpage

\section{\textcolor{orange!80!black}{第二课\quad 力——推和拉的秘密}}

\subsection*{力是什么？}

你推开门、提起书包、踢一脚足球——这些动作都有一个共同点：你在\textbf{对物体施加力}。力看不见，但它的效果人人都能感觉到。

\begin{definition}{力}
\textbf{力}是物体对物体的作用。力的作用是\textbf{相互的}——你打墙一拳，墙也同时打了你一拳（你的手疼就是证据）。
\end{definition}

力有两个重要的效果：

\textbf{效果一：改变物体的形状。}你用力捏橡皮泥，它变形了；你拉弹簧，它伸长了。这种变形有些是弹性的（力撤了能恢复，比如弹簧），有些是塑性的（力撤了也恢复不了，比如捏扁的易拉罐）。

\textbf{效果二：改变物体的运动状态。}运动的快慢变了（加速/减速/停下），或者运动的方向变了（拐弯），都叫"运动状态改变"。力是改变运动状态的原因——这句话非常重要，后面讲牛顿定律的时候会反复用到。

\subsection*{力的三要素}

描述一个力，你需要说清楚三件事：
\begin{enumerate}
  \item \textbf{大小}：力有多"大"。力的单位是\textbf{牛顿（N）}。1牛顿大约是你托起两个鸡蛋所用的力。
  \item \textbf{方向}：力朝哪边推（拉）。
  \item \textbf{作用点}：力作用在物体的哪个位置。
\end{enumerate}

在物理中，我们用一个带箭头的线段来表示力——\textbf{力的示意图}。线段的起点（或终点）是作用点，箭头指向力的方向，线段的长度按比例代表力的大小。画力的示意图是初中物理最重要的基本功之一。

\begin{figplaceholder}
此处插入力的三要素与力的示意图（见 figure/requirements/physics-mechanics1-ch1.txt 插图2）
\end{figplaceholder}

\subsection*{力的作用是相互的}

这可能是力学中最容易被忽略的规律——\textbf{物体间力的作用是相互的}。一对相互作用力\textbf{大小相等、方向相反、作用在同一条直线上，但作用在不同的物体上}。

举个例子：你站在地上，地球用重力拉你（向下），你的脚同时用同样大小的力压地面（向下）。而地面呢？地面用一个叫\textbf{支持力}的力把你托住（向上）。支持力和压力大小相等、方向相反，但它们不是一对相互作用力——它们是一对\textbf{平衡力}（作用在同一个物体——也就是你身上）。

\textcolor{gray}{\small 注：平衡力和相互作用力是初中物理最常考的辨析题。记住核心区别——平衡力作用在\textbf{同一物体}上；相互作用力作用在\textbf{不同物体}上。}

\begin{thinkbox}
划船的时候，桨向后划水，水给桨一个向前的力，船就往前走了。如果没有水——在太空中——你拼命挥桨也动不了。这就是力的相互性在起作用。想想看：火箭在太空中没有空气可以推，它是怎么飞的？（提示：燃料向下高速喷出，火箭受到向上的反冲力——和划船一个道理）
\end{thinkbox}

\begin{practice}
1. （判断）力的作用是相互的，所以两个力的大小一定相等。（\quad）\\
2. 托起两个鸡蛋所用的力大约是\_\_\_\_\_ N。\\
3. （选择）关于力的三要素，下列说法错误的是（\quad）\\
\quad A. 力的大小影响力的作用效果 \quad B. 力的方向影响力的作用效果\\
\quad C. 力的作用点影响力的作用效果 \quad D. 力的单位是千克\\
4. （填空）甲推乙，甲对乙的力和乙对甲的力是一对\_\_\_\_\_力，它们大小\_\_\_\_\_、方向\_\_\_\_\_、作用在\_\_\_\_\_物体上。\\
5. （选择）下列现象中说明力的作用是相互的是（\quad）\\
\quad A. 苹果从树上落下 \quad B. 船桨向后划水，船向前行 \quad C. 风吹树叶飘动 \quad D. 磁铁吸引铁钉\\
6. （填空）力可以改变物体的\_\_\_\_\_，也可以改变物体的\_\_\_\_\_。\\
7. 用绳子提水桶时，手受到的向下的拉力，施力物体是\_\_\_\_\_，受力物体是\_\_\_\_\_。\\
8. （判断）力的示意图中，线段的长度必须按比例代表力的大小。（\quad）
\end{practice}

\newpage

\section{\textcolor{orange!80!black}{第三课\quad 重力与弹力}}

\subsection*{重力——地球的"拥抱"}

如果你在自然科学小故事中读过牛顿与苹果的故事，你已经有直觉了——地球在"拉"着一切有质量的物体。这个拉力就是\textbf{重力}。

\begin{definition}{重力}
由于地球的吸引而使物体受到的力叫做\textbf{重力}，用 $G$ 表示。重力的方向总是\textbf{竖直向下}（指向地心）。
\end{definition}

重力的大小和物体的\textbf{质量}成正比：

\formula{G = mg}

其中 $m$ 是质量（单位：千克 kg），$g$ 是一个常数，约等于 $9.8\ \text{N/kg}$。在初中阶段，题目中通常取 $g = 10\ \text{N/kg}$ 方便计算。

\textbf{一个非常重要的辨析：质量和重量不是一回事！}质量是物体含有"物质多少"的量度——不管你在北京、在月球、还是在太空中，你的质量都不变。重量是重力的大小——你在月球上的重量只有地球上的大约六分之一（因为月球引力小），但你的质量没有变小。

\begin{example}
小明的质量是50 kg。他在地球上的重力是多少？如果他去了月球（$g_{\text{月}} \approx 1.6\ \text{N/kg}$），重力又变成多少？

\textbf{解：} 地球上：$G = mg = 50 \times 10 = 500\ \text{N}$。\\
月球上：$G = mg_{\text{月}} = 50 \times 1.6 = 80\ \text{N}$，不到地球上的六分之一！
\end{example}

\subsection*{弹力——恢复原状的力量}

你拉一条橡皮筋——它变长了。你放手——它缩回去了。这就是\textbf{弹性}。

\begin{definition}{弹力}
物体由于发生\textbf{弹性形变}而产生的力叫做\textbf{弹力}。弹力的方向与形变方向相反——你往右拉它，它往左"拽"你。
\end{definition}

弹簧测力计就是利用弹力原理制作的。在弹簧的弹性限度内，弹簧的伸长量与受到的拉力成正比——这个规律叫\textbf{胡克定律}。超过弹性限度，弹簧就废了——变形不可恢复。

生活中弹力无处不在：你坐的椅子（里面有弹簧或海绵）、你穿的球鞋（鞋底的弹性材料）、你用的橡皮——都在利用弹力缓冲或恢复。

\begin{thinkbox}
找一个弹簧测力计（或者橡皮筋+尺子代替），挂上不同重量的物体，记录伸长量。你会发现伸长量和重量基本成正比。这就是弹簧秤能称重的原因。但这个关系只有在一定范围内成立——如果你挂太重的东西把弹簧拉坏了，它就不准了。
\end{thinkbox}

\begin{practice}
1. （计算）一个质量为5kg的物体，在地球上的重力是多少？（$g = 10\ \text{N/kg}$）\\
2. （判断）宇航员在太空中质量变为零。（\quad）\\
3. （选择）关于质量和重力的关系，下列说法正确的是（\quad）\\
\quad A. 质量越大重力一定越大 \quad B. 质量与重力成正比，比值为 $g$\\
\quad C. 重力与质量成正比，比值为 $g$ \quad D. 重力就是质量\\
4. （填空）弹簧测力计的原理是：在\_\_\_\_\_限度内，弹簧的\_\_\_\_\_与受到的拉力成\_\_\_\_\_。\\
5. （计算）一个物体重力为120N，它的质量是多少？（$g=10\ \text{N/kg}$）\\
6. （选择）下列说法正确的是（\quad）\\
\quad A. 物体的质量越大，重力一定越大 \quad B. 1 kg = 10 N \quad C. 重力的方向总是垂直向下 \quad D. 宇航员在月球上没有质量\\
7. （填空）弹簧的劲度系数反映了弹簧的"\_\_\_\_\_"程度，劲度系数越大，弹簧越\_\_\_\_\_。\\
8. （判断）物体的重心一定在物体上。（\quad）
\end{practice}

\newpage

\section{\textcolor{orange!80!black}{第四课\quad 摩擦力——看不见的"阻力"}}

\subsection*{为什么地面上的东西自己不会一直滑？}

如果你在光滑的冰面上推一本书——书滑出去以后，速度越来越慢，最后停了。如果你在粗糙的水泥地上推同一本书——书滑了很短一段距离就停了。

这不是因为"书不想滑了"，而是\textbf{摩擦力}在起作用。

\begin{definition}{摩擦力}
两个互相接触的物体，当它们发生\textbf{相对运动}（或具有相对运动的趋势）时，在接触面上产生的一种阻碍相对运动的力，叫做\textbf{摩擦力}。
\end{definition}

\subsection*{三种摩擦力}

\textbf{静摩擦力：}两个物体接触但还没有相对滑动——比如你推一个很重的箱子，推不动。这时候箱子受到地面给的静摩擦力，大小等于你推的力，方向相反。静摩擦力有一个最大值——你力气够大的时候，箱子就开始动了。

\textbf{滑动摩擦力：}物体已经开始相对滑动了。滑动摩擦力的大小取决于两个因素：
\begin{itemize}
  \item 接触面的\textbf{粗糙程度}：越粗糙，摩擦力越大。
  \item 接触面之间的\textbf{压力大小}：压力越大，摩擦力越大。
\end{itemize}
\textbf{注意：}滑动摩擦力的大小\textbf{和接触面积的大小基本无关}。一块砖平着拉和立着拉，摩擦力几乎一样大——因为压力相同。

\textbf{滚动摩擦力：}一个物体在另一个物体表面滚动时产生的阻力。同样条件下，滚动摩擦远小于滑动摩擦——这就是为什么行李箱要装轮子、古人搬运巨石时在下面垫圆木。

\subsection*{增大和减小摩擦——都是学问}

有时候我们需要\textbf{增大摩擦}：鞋底花纹、轮胎花纹、体操运动员往手上涂镁粉（吸汗增加摩擦）、冬天路面结冰撒沙子。

有时候我们需要\textbf{减小摩擦}：门轴加润滑油、滑冰鞋的冰刀（压强融化冰形成水膜）、气垫船（在船底和地面之间形成气垫）、磁悬浮列车（用磁力让列车"飘"起来）。

\begin{figplaceholder}
此处插入三种摩擦力对比示意图：静摩擦/滑动摩擦/滚动摩擦（见 figure/requirements/physics-mechanics1-ch1.txt 插图3）
\end{figplaceholder}

\begin{example}
一个重100N的箱子放在水平地面上。如果用20N的水平推力推不动，摩擦力多大？如果用50N的力刚好推动，最大静摩擦力多大？如果推动后只需40N的力就能保持匀速运动，滑动摩擦力多大？

\textbf{解：}\\
推不动：箱子静止，二力平衡 $\rightarrow$ 静摩擦力 = 20N（等于推力）。\\
刚好推动：最大静摩擦力 = 50N。\\
匀速滑动：匀速运动也是平衡状态 $\rightarrow$ 滑动摩擦力 = 40N（等于维持匀速的推力）。
\end{example}

\begin{thinkbox}
假如世界上突然没有摩擦力了——会怎样？你走路会像在冰面上一样滑倒（脚蹬地没有摩擦力就无法前进）；钉子会从墙上滑出来；所有绳结都会自己松开；汽车无法起步也无法刹车……摩擦力虽然叫"阻力"，但我们的生活一刻也离不开它。
\end{thinkbox}

\begin{practice}
1. （判断）摩擦力总是阻碍物体的运动。（\quad）\\
2. （选择）下列措施中属于增大摩擦的是（\quad）\\
\quad A. 门轴加润滑油 \quad B. 鞋底做花纹 \quad C. 气垫船形成气垫 \quad D. 行李箱装轮子\\
3. （填空）推一个重100N的箱子，用30N力推不动——此时摩擦力为\_\_\_\_\_N；用45N力匀速推动——滑动摩擦力为\_\_\_\_\_N。\\
4. 滚动摩擦远\_\_\_\_\_于滑动摩擦（填"大"或"小"），所以搬运重物时常在下面垫\_\_\_\_\_。\\
5. （计算）一个重200N的箱子放在水平地面上，用80N的水平力推但没有推动，摩擦力为\_\_\_\_\_N。\\
6. （选择）下列做法中为了减小摩擦的是（\quad）\\
\quad A. 鞋底做花纹 \quad B. 自行车刹车时捏紧闸 \quad C. 给链条加润滑油 \quad D. 瓶盖刻有花纹\\
7. （填空）滑动摩擦力的大小与\_\_\_\_\_和\_\_\_\_\_有关。\\
8. （判断）物体的运动速度越大，滑动摩擦力越大。（\quad）
\end{practice}

\vspace{0.5cm}

\newpage

\section{\textcolor{orange!80!black}{第五课\quad 二力平衡与惯性}}

\subsection*{静止和匀速——平衡状态}

你面前有一本书静静地躺在桌子上。它受不受力？当然受——地球拉它（重力，向下），桌子托它（支持力，向上）。但书纹丝不动。为什么？因为这两个力\textbf{互相抵消了}——它们大小相等、方向相反，作用在同一个物体上。

\begin{definition}{二力平衡}
一个物体在两个力的作用下，如果保持\textbf{静止}或\textbf{匀速直线运动}状态，我们就说这两个力是\textbf{平衡力}，或者说物体处于\textbf{平衡状态}。
\end{definition}

二力平衡的条件有四个（缺一不可）：
\begin{enumerate}
  \item 作用在\textbf{同一个物体}上
  \item 大小\textbf{相等}
  \item 方向\textbf{相反}
  \item 作用在\textbf{同一条直线}上
\end{enumerate}

\textbf{重要辨析——平衡力 vs 相互作用力：}前面讲过，相互作用力也满足"大小相等、方向相反、作用在同一直线上"，但——相互作用力是作用在\textbf{两个不同物体}上（A推B，B推A）；而平衡力是作用在\textbf{同一个物体}上。这是考试中最经典的辨析题，请读三遍。

\subsection*{牛顿第一定律——惯性定律}

如果我们把桌面上那本书往外推一下——它滑出去了。然后呢？它越滑越慢，最后停了。为什么停？因为有\textbf{摩擦力}在阻挠它。

可是，如果在绝对光滑的表面上推这本书——它会不会永远滑下去？

牛顿的回答是：\textbf{会的。}

\begin{definition}{牛顿第一定律（惯性定律）}
一切物体在没有受到外力作用的时候，总保持\textbf{静止状态}或\textbf{匀速直线运动状态}。
\end{definition}

这个定律告诉我们两件事：
\begin{enumerate}
  \item \textbf{力不是维持运动的原因}——不受力的物体也能一直运动（匀速直线运动）。
  \item \textbf{力是改变运动状态的原因}——物体之所以停下来、加速、减速、拐弯，都是因为受了力。
\end{enumerate}

这两句话是物理学史上最重要的观念转变之一。古希腊的亚里士多德认为"物体不推就不动"——几千年人们都信他。牛顿说：不对，不推也能动；推了反而会改变运动状态。

\subsection*{惯性}

牛顿第一定律又叫做\textbf{惯性定律}。惯性是物体自身的属性——\textbf{一切物体都有保持原来运动状态不变的性质}。

惯性的大小只和\textbf{质量}有关——质量越大，惯性越大，越难改变它的运动状态。你推一辆自行车很容易，推一辆汽车就几乎推不动——不是因为汽车受了更大的摩擦力，而是它的惯性大得多。同样道理，高速行驶的大卡车比小轿车难刹车得多——你刹车踩死了，巨大的惯性还在"拽"着它往前冲。

\textbf{生活中的惯性：}急刹车时人往前倾（身体保持原来向前运动的状态）；抖衣服可以把灰尘抖掉（灰尘由于惯性留在原地，衣服跑了）；锤头松了，把锤柄往地上撞几下就紧了（锤头保持向下运动的惯性）。

\begin{example}
一辆汽车正在匀速直线行驶，突然左转弯。车里的乘客会往哪边倒？为什么？

\textbf{解：}乘客会往\textbf{右侧}倒。因为乘客的身体具有惯性，倾向于保持原来的向前运动方向。汽车向左转，乘客的上半身来不及跟着左转，在车里看起来就是向右侧倾倒。
\end{example}

\begin{thinkbox}
在太空中（没有重力、没有摩擦力），你扔出一个球——它会怎样运动？答案是：永远沿直线匀速飞下去，直到被什么挡到。这就是牛顿第一定律最"纯粹"的体现。地球上的物体之所以停，都是因为摩擦力、空气阻力这些"外力"在捣乱。
\end{thinkbox}

\begin{practice}
1. （判断）物体不受力时，一定处于静止状态。（\quad）\\
2. （选择）关于惯性，下列说法正确的是（\quad）\\
\quad A. 物体运动越快，惯性越大 \quad B. 物体静止时没有惯性\\
\quad C. 惯性大小只与质量有关 \quad D. 太空中的物体没有惯性\\
3. 汽车急刹车时，乘客往前倾倒，是因为乘客具有\_\_\_\_\_。\\
4. （辨析）一本书放在桌面上静止。书的重力和桌面对书的支持力是一对\_\_\_\_\_力（填"平衡"或"相互作用"）；书对桌面的压力和桌面对书的支持力是一对\_\_\_\_\_力。\\
5. （选择）下列物体受到平衡力作用的是（\quad）\\
\quad A. 从斜面加速下滑的物体 \quad B. 水平路面上匀速行驶的汽车 \quad C. 刚从枪膛射出的子弹 \quad D. 绕地球匀速转动的卫星\\
6. （填空）汽车在平直公路上匀速行驶时，水平方向上受到的\_\_\_\_\_力和\_\_\_\_\_力是一对平衡力。\\
7. 锤头松了，把锤柄往地上撞几下就紧了，是利用了\_\_\_\_\_的惯性。\\
8. （判断）物体只有在不受力时才具有惯性。（\quad）
\end{practice}

\newpage

\section{\textcolor{orange!80!black}{第六课\quad 压强——为什么针尖能刺穿}}

\subsection*{同样的力，不同的效果}

你用同样的力气(比如50N)，用手掌压一块木板——什么也没发生。换成一根针，用同样的力气扎下去——针尖穿透了木板。

为什么同样的力，效果天差地别？因为\textbf{力分散的面积不一样}。

\begin{definition}{压强}
物体单位面积上受到的压力叫做\textbf{压强}，用 $p$ 表示。
\end{definition}

\formula{p = \frac{F}{S}}

其中 $F$ 是压力（单位：N），$S$ 是受力面积（单位：$\text{m}^2$）。压强的单位是\textbf{帕斯卡（Pa）}，$1\ \text{Pa} = 1\ \text{N/m}^2$。

这个公式告诉我们：\textbf{压强和压力成正比，和受力面积成反比。}

\begin{itemize}
  \item 压力相同时——受力面积越小，压强越大（所以针能刺穿、刀刃能切断）。
  \item 受力面积相同时——压力越大，压强越大（同样的鞋底，大人踩你比小孩踩你疼）。
\end{itemize}

\subsection*{增大压强和减小压强——生活的智慧}

\textbf{增大压强的方法：}增大压力，或减小受力面积。钉子、针、刀刃、注射器针头——越尖锐，受力面越小，压强越大。

\textbf{减小压强的方法：}减小压力，或增大受力面积。坦克的履带（极大的接地面积——沼泽地上都不会陷下去）、书包的宽背带（增大和肩膀的接触面积以减小压强——背同样的书更舒服）、铁轨下面的枕木（把火车的重量分散到更大的地面上）。

\begin{example}
一本物理书平放在桌面上，质量0.3 kg，封面长26cm、宽18cm。书对桌面的压强是多少？（$g = 10\ \text{N/kg}$）

\textbf{解：} 书对桌面的压力 $F = mg = 0.3 \times 10 = 3\ \text{N}$。\\
受力面积 $S = 0.26 \times 0.18 = 0.0468\ \text{m}^2$。\\
$p = \dfrac{F}{S} = \dfrac{3}{0.0468} \approx 64.1\ \text{Pa}$。

这个压强非常小——大概就相当于几片纸压在手上。但如果把这本书竖起来、用书脊接触桌面（书脊只有约1cm宽），压强就会暴增几十倍。
\end{example}

\begin{figplaceholder}
此处插入压强原理示意图：手掌 vs 针尖（同样力不同面积）+ 生活实例（坦克履带/书包带/刀刃）（见 figure/requirements/physics-mechanics1-ch2.txt 插图1）
\end{figplaceholder}

\begin{practice}
1. （计算）一个重500N的人，双脚站立时与地面的接触面积约$0.04\ \text{m}^2$。他对地面的压强是多少？\\
2. （选择）下列实例中，属于减小压强的是（\quad）\\
\quad A. 针做得尖锐 \quad B. 菜刀刀刃很薄 \quad C. 书包带做得很宽 \quad D. 图钉的钉尖很尖\\
3. （填空）压强公式 $p = F/S$ 中，$F$ 的单位是\_\_\_\_\_，$S$ 的单位是\_\_\_\_\_。\\
4. 坦克安装宽大履带是为了\_\_\_\_\_压强；刀刃磨薄是为了\_\_\_\_\_压强（填"增大"或"减小"）。\\
5. （计算）一辆坦克质量20t，每条履带与地面接触面积2m$^2$。坦克对地面的压强多大？（$g=10\ \text{N/kg}$）\\
6. （选择）下列实例中属于增大压强的是（\quad）\\
\quad A. 铁轨铺在枕木上 \quad B. 骆驼宽大脚掌 \quad C. 啄木鸟细长的喙 \quad D. 推土机宽大履带\\
7. （填空）压强是表示\_\_\_\_\_的物理量，国际单位是\_\_\_\_\_。\\
8. 铅笔两端顶在手指上，\_\_\_\_\_端手指感觉更疼，因为\_\_\_\_\_。
\end{practice}

\newpage

\section{\textcolor{orange!80!black}{第七课\quad 液体压强与大气压强}}

\subsection*{液体也有压强——而且四面八方都有}

固体压强的方向是向下的（压力方向），但\textbf{液体压强}不同——液体在\textbf{各个方向}都有压强。你把一个塑料袋套在手上放进水里，你会发现袋子从各个方向紧贴你的手——上、下、左、右都有压力。

液体压强的大小取决于两个因素：液体的\textbf{密度}和\textbf{深度}。

\formula{p = \rho g h}

其中 $\rho$ 是液体的密度（$\text{kg/m}^3$），$g = 10\ \text{N/kg}$，$h$ 是从液面到所求点的垂直深度（m）。

这个公式有几个重要推论：
\begin{itemize}
  \item 同一液体中，深度越大，压强越大——所以潜水越深，身体承受的水压越大。深海潜水需要特制的潜水服和潜艇。
  \item 同一深度，液体向各个方向的压强\textbf{大小相等}。
  \item 不同液体在同一深度，密度越大的液体产生的压强越大——所以同样潜水一米，在海里比在淡水湖里承受的压强略大（海水密度约$1.03 \times 10^3\ \text{kg/m}^3$）。
\end{itemize}

\textbf{连通器原理：}底部互相连通的容器（连通器）里，如果只有同一种液体，在液体不流动时，各容器中的液面\textbf{总是相平的}。茶壶的壶嘴和壶身、锅炉的水位计、三峡大坝的船闸——都是连通器的应用。

\begin{figplaceholder}
此处插入液体压强示意图：深度与压强的关系+连通器原理（见 figure/requirements/physics-mechanics1-ch2.txt 插图2）
\end{figplaceholder}

\subsection*{大气压强——整片天空压着你}

你头顶上有一百多公里厚的大气层。空气虽然轻，但\textbf{积少成多}——在海平面上，大气压强约为 $1.013 \times 10^5\ \text{Pa}$。这意味着每平方米的地面上压着大约\textbf{10吨重}的空气！

那为什么你没被压扁？因为你身体内部也有气体（肺里的空气、体液里溶解的气体），内外压强基本平衡。

\textbf{托里拆利实验：}1643年，意大利科学家托里拆利用一根装满水银的玻璃管倒插在水银槽里——管里的水银柱下降到了一定高度（约760mm）就不下降了。这个76厘米高的水银柱就是被大气压强"托"住的。这个实验第一次精确测量了大气压强的数值。

\textbf{大气压强随高度变化：}海拔越高，头顶的空气越薄，大气压越低。这就是为什么坐飞机时耳朵会胀痛（中耳内的气压来不及平衡外界降低的气压），以及为什么在高原烧水不到100℃就沸腾了（气压越低，水的沸点越低）。

\begin{example}
一个标准大气压 $p_0 = 1.0 \times 10^5\ \text{Pa}$。它相当于多高的水柱产生的压强？（$\rho_{\text{水}} = 1.0 \times 10^3\ \text{kg/m}^3$，$g = 10\ \text{N/kg}$）

\textbf{解：} $p_0 = \rho g h \quad \rightarrow \quad h = \dfrac{p_0}{\rho g} = \dfrac{1.0 \times 10^5}{1.0 \times 10^3 \times 10} = 10\ \text{m}$。

一个标准大气压相当于10米高的水柱产生的压强！这意味着如果你用一根超过10米长的吸管吸水，不管你多用力也吸不上来——不是因为你不够努力，是大气压只能把水"推"到10米高。
\end{example}

\begin{thinkbox}
拿一个玻璃杯装满水，用一张纸片盖住杯口，按住纸片快速把杯子倒过来——然后松手。纸片居然不掉，水也不流出来！是什么托住了这杯水？大气压强。托里拆利用76厘米水银柱证明的事，你用一杯水就能亲手体验。
\end{thinkbox}

\begin{practice}
1. （计算）一个游泳池水深3米处，水的压强多大？（$\rho = 1.0\times 10^3\ \text{kg/m}^3$，$g=10\ \text{N/kg}$）\\
2. （判断）连通器中各容器的液面一定相平。（\quad）\\
3. 托里拆利实验中，玻璃管倾斜时，管内水银柱的竖直高度\_\_\_\_\_（填"变大""变小"或"不变"）。\\
4. （选择）高原上烧水不到100℃就沸腾了，是因为（\quad）\\
\quad A. 高原上温度低 \quad B. 高原上气压低 \quad C. 高原上氧气少 \quad D. 高原上火不够大\\
5. （填空）一个标准大气压约等于\_\_\_\_\_ Pa。\\
6. （计算）潜水员在水下20m深处承受的压强多大？（$\rho_{\text{水}}=1.0\times10^3\ \text{kg/m}^3$，$g=10\ \text{N/kg}$）\\
7. （选择）下列应用利用连通器原理的是（\quad）\\
\quad A. 温度计 \quad B. 茶壶 \quad C. 吸盘挂钩 \quad D. 注射器\\
8. （填空）大气压随海拔高度增加而\_\_\_\_\_（增大/减小），水的沸点随气压降低而\_\_\_\_\_。\\
9. 用吸管喝饮料是利用了\_\_\_\_\_将饮料压入嘴中。
\end{practice}

\newpage

\section{\textcolor{orange!80!black}{第八课\quad 浮力——阿基米德的"升级版"}}

\subsection*{从浴缸到公式}

如果你读过自然科学小故事，你一定记得阿基米德光着身子从浴缸跳出来大喊"尤里卡"的故事。他发现：\textbf{物体浸入水中时，排开的水的重量，等于物体受到的浮力。}

现在我们要把这个发现写成精准的物理公式：

\formula{F_{\text{浮}}} = \rho_{\text{液}} g V_{\text{排}}}

其中 $\rho_{\text{液}}$ 是液体的密度（$\text{kg/m}^3$），$g = 10\ \text{N/kg}$，$V_{\text{排}}$ 是物体排开液体的体积（$\text{m}^3$）。

这个公式叫\textbf{阿基米德原理}。注意两个要点：
\begin{enumerate}
  \item $V_{\text{排}}$ 不一定等于物体的总体积——如果物体只浸入了一部分（比如浮在水面上），$V_{\text{排}}$ 只是浸入的那部分体积。
  \item 浮力的大小只取决于液体的密度和排开的体积——和物体自身的密度、形状、在水下的深度都\textbf{没有直接关系}。
\end{enumerate}

\subsection*{物体的沉浮条件}

为什么有的东西浮、有的东西沉？不是看物体"多重"——是看\textbf{重力和浮力谁大}。

\begin{center}
\begin{tabular}{|c|c|c|}
\hline
\textbf{状态} & \textbf{力的关系} & \textbf{密度的关系} \\
\hline
下沉（沉底） & $G > F_{\text{浮}}$ & $\rho_{\text{物}} > \rho_{\text{液}}$ \\
悬浮（停在水中）& $G = F_{\text{浮}}$ & $\rho_{\text{物}} = \rho_{\text{液}}$ \\
上浮（最终漂浮）& $G < F_{\text{浮}}$（上浮过程）& $\rho_{\text{物}} < \rho_{\text{液}}$ \\
漂浮（停在水面）& $G = F_{\text{浮}}$ & $\rho_{\text{物}} < \rho_{\text{液}}$ \\
\hline
\end{tabular}
\end{center}

\textbf{注意：悬浮和漂浮不一样。}悬浮是物体\textbf{完全浸没}在水中，停在任何深度；漂浮是物体\textbf{部分露出}水面。

\textbf{铁船为什么能浮？}铁块的密度（$7.9 \times 10^3\ \text{kg/m}^3$）远大于水——一块实心铁肯定沉底。但把铁做成空心的船——船壳是铁，但船内部是空的，整体\textbf{平均密度}小于水。船的排水量就是它排开的水的重量——万吨巨轮排开巨量的水，产生巨量的浮力。

\textbf{潜水艇的浮沉原理：}潜水艇通过往水舱里注水（增加自身重力）或排出水（减少自身重力）来控制沉浮——在这个过程中，$V_{\text{排}}$ 不变（完全浸没时），靠改变自身重力来达到 $G > F_{\text{浮}}$ 或 $G = F_{\text{浮}}$。

\begin{example}
一个体积为 $0.001\ \text{m}^3$（即1升）的铁块（$\rho_{\text{铁}} = 7.9 \times 10^3\ \text{kg/m}^3$），完全浸没在水中。它受到的浮力多大？它会上浮还是下沉？

\textbf{解：} 完全浸没时 $V_{\text{排}} = V_{\text{物}} = 0.001\ \text{m}^3$。\\
$F_{\text{浮}} = \rho_{\text{水}} g V_{\text{排}} = 1000 \times 10 \times 0.001 = 10\ \text{N}$。\\
铁块的重力 $G = mg = \rho_{\text{铁}} V g = 7900 \times 0.001 \times 10 = 79\ \text{N}$。\\
$G(79\ \text{N}) > F_{\text{浮}}(10\ \text{N})$ ——下沉。
\end{example}

\begin{practice}
1. （计算）一个体积为$0.5\ \text{m}^3$的木块漂浮在水面上，浸入水中的体积为$0.3\ \text{m}^3$。它受到的浮力多大？（$\rho_{\text{水}}=1000\ \text{kg/m}^3$，$g=10\ \text{N/kg}$）\\
2. （判断）物体在水中下沉是因为没有受到浮力。（\quad）\\
3. （选择）一艘轮船从长江驶入大海，它受到的浮力（\quad）\\
\quad A. 变大 \quad B. 变小 \quad C. 不变 \quad D. 无法判断\\
4. （填空）当 $G$\_\_\_\_\_$F_{\text{浮}}$ 时物体下沉；当物体漂浮时，$G$\_\_\_\_\_$F_{\text{浮}}$（填 $>$、$<$ 或 $=$）。\\
5. （计算）一物体在空气中重10N，浸没在水中时弹簧测力计示数为6N。它受到的浮力多大？\\
6. （选择）轮船从海中驶入江中，受到的浮力（\quad）\\
\quad A. 变大 \quad B. 变小 \quad C. 不变 \quad D. 无法判断\\
7. （填空）潜水艇是通过改变\_\_\_\_\_来实现上浮和下潜的。\\
8. 同一鸡蛋在盐水中漂浮、在水中沉底，浮力：$F_{\text{盐水}}$\_\_\_\_\_$F_{\text{水}}$（填$>$、$<$或$=$）。\\
9. （判断）物体排开液体的体积越大，受到的浮力一定越大。（\quad）
\end{practice}

\newpage

\section{\textcolor{orange!80!black}{第九课\quad 简单机械——省力的智慧}}

\subsection*{杠杆——阿基米德的"撬动地球"}

阿基米德说过一句著名的话："给我一个支点，我能撬动地球。"他说的就是\textbf{杠杆}。

\begin{definition}{杠杆}
一根在力的作用下能绕\textbf{固定点}转动的硬棒叫做\textbf{杠杆}。
\end{definition}

杠杆有五要素：
\begin{enumerate}
  \item \textbf{支点 $O$}：杠杆绕着转动的固定点
  \item \textbf{动力 $F_1$}：使杠杆转动的力
  \item \textbf{阻力 $F_2$}：阻碍杠杆转动的力
  \item \textbf{动力臂 $l_1$}：支点到动力作用线的垂直距离
  \item \textbf{阻力臂 $l_2$}：支点到阻力作用线的垂直距离
\end{enumerate}

\textbf{注意：力臂是"支点到力的作用线的垂直距离"，不是支点到作用点的距离。}画力臂的时候，要从支点向力的作用线画垂线——这是很多人容易犯的错误。

\begin{figplaceholder}
此处插入杠杆五要素示意图：标注支点/动力/阻力/动力臂/阻力臂（见 figure/requirements/physics-mechanics1-ch2.txt 插图3）
\end{figplaceholder}

\subsection*{杠杆的平衡条件}

当杠杆在动力和阻力的作用下静止时，我们说杠杆平衡了。杠杆平衡的条件是：

\formula{F_1 \cdot l_1 = F_2 \cdot l_2 \qquad \text{即} \qquad \text{动力} \times \text{动力臂} = \text{阻力} \times \text{阻力臂}}

根据动力臂和阻力臂的长短关系，杠杆分三类：
\begin{itemize}
  \item \textbf{省力杠杆}（$l_1 > l_2$）：省力，但费距离。撬棍、开瓶器、手推车。
  \item \textbf{费力杠杆}（$l_1 < l_2$）：费力，但省距离。筷子、镊子、钓鱼竿。为什么筷子费力还要用？因为用筷子可以精确控制食物——手指移动一小段，筷子尖就移动一段合适的大距离。
  \item \textbf{等臂杠杆}（$l_1 = l_2$）：不省力也不费力。天平。
\end{itemize}

\subsection*{滑轮——改变力的方向}

\textbf{定滑轮：}轴固定不动的滑轮。它不省力，但可以\textbf{改变力的方向}——向下拉绳子，重物向上升。旗杆顶上的滑轮就是定滑轮。

\textbf{动滑轮：}轴随重物一起移动的滑轮。它省力——理想情况下可以省一半力（$F = G/2$），但\textbf{不能改变力的方向}，而且绳子要多拉一倍距离。

\textbf{滑轮组：}把定滑轮和动滑轮组合在一起——既省力，又能改变力的方向。滑轮组是实际应用最多的方案：塔吊、电梯、起重机都用滑轮组。

\begin{example}
用滑轮组将600N的重物匀速提升。滑轮组由3段绳子承担物重。忽略摩擦和动滑轮重，拉力多大？绳子要拉多少米才能将重物提升2米？

\textbf{解：} 3段绳子承担物重 $\rightarrow$ $F = \dfrac{G}{n} = \dfrac{600}{3} = 200\ \text{N}$。\\
绳子自由端移动距离 $s = n \cdot h = 3 \times 2 = 6\ \text{m}$。\\
省力不省功——力省了，路程翻了倍。
\end{example}

\begin{practice}
1. （计算）一根杠杆动力臂$1.5\ \text{m}$，阻力臂$0.3\ \text{m}$。要撬动$600\ \text{N}$的重物，至少需要多大的动力？\\
2. （选择）下列工具中属于费力杠杆的是（\quad）\\
\quad A. 开瓶器 \quad B. 筷子 \quad C. 撬棍 \quad D. 羊角锤\\
3. （填空）定滑轮\_\_\_\_\_力，但可以改变\_\_\_\_\_；动滑轮可以\_\_\_\_\_力，但不能改变\_\_\_\_\_。\\
4. 使用滑轮组时，承担物重的绳子段数越多，越\_\_\_\_\_力（填"省"或"费"）。\\
5. （选择）下列工具中属于省力杠杆的是（\quad）\\
\quad A. 镊子 \quad B. 钓鱼竿 \quad C. 钢丝钳 \quad D. 筷子\\
6. （填空）杠杆的平衡条件表达式是\_\_\_\_\_。\\
7. （计算）用滑轮组将800N重物提升2m，由4段绳子承担物重。拉力至少多大？（不计摩擦和动滑轮重）\\
8. （判断）使用滑轮组既能省力又能省距离。（\quad）\\
9. （填空）定滑轮的本质是\_\_\_\_\_杠杆，动滑轮的本质是\_\_\_\_\_杠杆。
\end{practice}

\newpage

\section{\textcolor{orange!80!black}{第十课\quad 功与机械效率}}

\subsection*{物理中的"做功"}

在日常生活中，"做功"就是"干活"。在物理中，\textbf{做功}有更严格的含义：

\begin{definition}{功}
如果一个力作用在物体上，并且物体在\textbf{这个力的方向上}移动了一段距离，我们就说这个力对物体做了功。
\end{definition}

\formula{W = Fs}

其中 $W$ 是功（单位：焦耳 J），$F$ 是力（N），$s$ 是物体在力的方向上移动的距离（m）。$1\ \text{J} = 1\ \text{N} \cdot \text{m}$。

\textbf{做功的两个必要因素：}
\begin{enumerate}
  \item 有力作用在物体上
  \item 物体在力的方向上移动了距离
\end{enumerate}
两个条件缺一不可。你拼命推一堵墙，墙纹丝不动——你做了功吗？在物理意义上，\textbf{没有}。因为 $s = 0$，$W = 0$。你消耗了体力（化学能），但没有对墙做机械功。

\subsection*{功率——做功的快慢}

两个人搬同样的砖到同样高度——做的功一样多。但一个人用了1分钟，另一个人用了10分钟——他们做功的快慢不同。

\begin{definition}{功率}
\textbf{功率}是描述做功快慢的物理量，等于单位时间内做的功。
\end{definition}

\formula{P = \frac{W}{t}}

功率的单位是\textbf{瓦特（W）}，$1\ \text{W} = 1\ \text{J/s}$。如果你读过自然科学小故事，你一定记得——这个单位就是用小瓦特的名字命名的。你家灯泡上"40W"就是40瓦特。

\subsection*{机械效率——"划不划算"}

任何机械在工作时都有摩擦和其他损耗。我们用机械做的"有用功"总是小于我们对机械输入的"总功"。\textbf{机械效率}衡量一台机械"划算"的程度：

\formula{\eta = \frac{W_{\text{有用}}}{W_{\text{总}}} \times 100\%}

机械效率永远\textbf{小于100\%}（现实中没有不损耗的永动机）。提高机械效率的常见方法：减小摩擦（润滑）、减轻机械自重（用轻质材料）。

\begin{example}
用滑轮组将200N的重物提升5米。拉力做的总功为1500J。有用功多大？机械效率是多少？

\textbf{解：} 有用功 = $Gh = 200 \times 5 = 1000\ \text{J}$。\\
$\eta = \dfrac{W_{\text{有用}}}{W_{\text{总}}} \times 100\% = \dfrac{1000}{1500} \times 100\% \approx 66.7\%$。\\
剩下33.3\%的能量去哪了？变成了克服摩擦和滑轮自重产生的热量——你没有浪费它，它只是"跑"到别的地方去了。
\end{example}

\begin{thinkbox}
物理中的"功"和生活中的"累"不是一回事。你举着书包站在公交车上——手酸得很（肌肉持续收缩消耗能量），但从物理意义上说，你没有对书包做机械功（书包在竖直方向没有移动距离 $s=0$）。你消耗的能量全变成了肌肉里热量的释放。
\end{thinkbox}

\begin{practice}
1. （计算）一个人用50N的水平推力推箱子前进了10m。他对箱子做了多少功？\\
2. （判断）做功的两个必要因素是\_\_\_\_\_和\_\_\_\_\_。\\
3. （计算）一台机器5秒内做功1000J，功率是多少瓦？\\
4. （选择）关于机械效率，下列说法正确的是（\quad）\\
\quad A. 机械效率可以等于100\% \quad B. 机械效率越大，机械越省力\\
\quad C. 机械效率是有用功与总功的比值 \quad D. 功率越大，机械效率越高\\
5. （填空）提高机械效率的方法有\_\_\_\_\_摩擦和减轻\_\_\_\_\_。\\
6. （计算）小明用100N的力将重20N的物体匀速提升3m。他对物体做了多少功？\\
7. （选择）下列情况中人对物体做功的是（\quad）\\
\quad A. 举着杠铃静止不动 \quad B. 提着水桶在水平路面匀速前进 \quad C. 扛着米袋上楼梯 \quad D. 推石头没推动\\
8. （填空）功率是表示\_\_\_\_\_的物理量，一台机器功率为1kW表示1秒内做功\_\_\_\_\_J。\\
9. （判断）机械效率越高的机械一定越省力。（\quad）
\end{practice}

\vspace{0.5cm}

\begin{center}
\rule{0.7\textwidth}{1pt}

\vspace{0.4cm}
{\large\textcolor{blue!70!black}{\textbf{—— 物理1·力学入门 完 ——}}}

\vspace{0.3cm}
\textcolor{gray}{\small 下一本预告：物理2·声光热——从声音的振动到透镜的魔法}
\end{center}

\newpage
\section*{\textcolor{defblue}{参考答案}}

\textbf{第一课 练一练：}
1. ×（以车窗为参照物，小明静止）\quad
2. 20 m/s，合72 km/h\quad
3. C\quad
4. 静止\quad
5. $t = \dfrac{s}{v} = \dfrac{120}{300} = 0.4\ \text{h} = 24\ \text{min}$\quad
6. $v = s/t = 1200/300 = 4\ \text{m/s}$，合14.4 km/h\quad
7. $\surd$\quad
8. C\quad
9. 运动快慢，$v=s/t$

\textbf{第二课 练一练：}
1. ×（不一定——"两个力"不明确时结论不成立）\quad
2. 1 N\quad
3. D（力的单位是牛顿）\quad
4. 相互作用力，相等，相反，不同\quad
5. B\quad
6. 形状，运动状态\quad
7. 绳子，手\quad
8. $\surd$

\textbf{第三课 练一练：}
1. $G = mg = 5\times10 = 50\ \text{N}$\quad
2. ×（质量不变）\quad
3. C\quad
4. 弹性，伸长量，正比\quad
5. $m = G/g = 120/10 = 12\ \text{kg}$\quad
6. A\quad
7. 硬，硬\quad
8. ×（重心不一定在物体上，如圆环）

\textbf{第四课 练一练：}
1. ×（摩擦力阻碍的是相对运动）\quad
2. B\quad
3. 30 N，45 N\quad
4. 小，圆木（或滚木）\quad
5. 80 N\quad
6. C\quad
7. 压力大小，接触面粗糙程度\quad
8. ×（滑动摩擦力与速度无关）

\textbf{第五课 练一练：}
1. ×（也可能做匀速直线运动）\quad
2. C\quad
3. 惯性\quad
4. 平衡力，相互作用力\quad
5. B\quad
6. 牵引力，摩擦力\quad
7. 锤头\quad
8. ×（一切物体任何时候都有惯性）

\textbf{第六课 练一练：}
1. $p = \dfrac{F}{S} = \dfrac{500}{0.04} = 12500\ \text{Pa}$\quad
2. C\quad
3. N，m$^2$\quad
4. 减小，增大\quad
5. $F=mg=2\times10^4\ \text{N}$，$p=F/(2S)=2\times10^4/4=5\times10^3\ \text{Pa}$\quad
6. C\quad
7. 压力作用效果，帕斯卡（Pa）\quad
8. 尖，受力面积小压强大

\textbf{第七课 练一练：}
1. $p = \rho gh = 1000\times10\times3 = 3\times10^4\ \text{Pa}$\quad
2. ×（必须是同种液体且不流动时相平）\quad
3. 不变\quad
4. B\quad
5. $1.013\times10^5$\quad
6. $p=\rho gh=1000\times10\times20=2\times10^5\ \text{Pa}$\quad
7. B\quad
8. 减小，降低\quad
9. 大气压强

\textbf{第八课 练一练：}
1. $F_{\text{浮}} = \rho_{\text{水}}gV_{\text{排}} = 1000\times10\times0.3 = 3000\ \text{N}$\quad
2. ×（下沉是因为重力大于浮力）\quad
3. C（漂浮时浮力始终等于重力）\quad
4. $>$，$=$\quad
5. $F_{\text{浮}} = 10-6 = 4\ \text{N}$\quad
6. C（漂浮时$F_{\text{浮}}=G$不变）\quad
7. 自身重力\quad
8. $>$\quad
9. ×（还与液体密度有关）

\textbf{第九课 练一练：}
1. $F_1 = \dfrac{F_2l_2}{l_1} = \dfrac{600\times0.3}{1.5} = 120\ \text{N}$\quad
2. B\quad
3. 不省，力的方向；省，力的方向\quad
4. 省\quad
5. C\quad
6. $F_1l_1 = F_2l_2$\quad
7. $F = G/n = 800/4 = 200\ \text{N}$\quad
8. ×（省力不省功，费距离）\quad
9. 等臂，省力

\textbf{第十课 练一练：}
1. $W = Fs = 50\times10 = 500\ \text{J}$\quad
2. 力，物体在力的方向上移动的距离\quad
3. $P = \dfrac{W}{t} = \dfrac{1000}{5} = 200\ \text{W}$\quad
4. C\quad
5. 减小，机械自重\quad
6. $W = Fs = 100\times3 = 300\ \text{J}$\quad
7. C\quad
8. 做功快慢，$P=W/t$\quad
9. 1000\quad
10. ×（机械效率与是否省力无关）

\end{document}

